\section{Text Examples} % Sections are added in order to organize your presentation into discrete blocks, all sections and subsections are automatically output to the table of contents as an overview of the talk but NOT output in the presentation as separate slides

%------------------------------------------------

\subsection{Paragraphs and Lists}

\begin{frame}
    \frametitle{Paragraphs of Text}

    Sed iaculis \alert{dapibus gravida}. Morbi sed tortor erat, nec interdum arcu. Sed id lorem lectus. Quisque viverra augue id sem ornare non aliquam nibh tristique. Aenean in ligula nisl. Nulla sed tellus ipsum. Donec vestibulum ligula non lorem vulputate fermentum accumsan neque mollis.

    \bigskip % Vertical whitespace

    % Quote example
    \begin{quote}
        Sed diam enim, sagittis nec condimentum sit amet, ullamcorper sit amet libero. Aliquam vel dui orci, a porta odio.\\
        --- Someone, somewhere\ldots
    \end{quote}

    \bigskip % Vertical whitespace

    Nullam id suscipit ipsum. Aenean lobortis commodo sem, ut commodo leo gravida vitae. Pellentesque vehicula ante iaculis arcu pretium rutrum eget sit amet purus. Integer ornare nulla quis neque ultrices lobortis.
\end{frame}

%------------------------------------------------

\begin{frame}
    \frametitle{Lists}
    \framesubtitle{Bullet Points and Numbered Lists} % Optional subtitle

    \begin{itemize}
        \item Lorem ipsum dolor sit amet, consectetur adipiscing elit
        \item Aliquam blandit faucibus nisi, sit amet dapibus enim tempus
              \begin{itemize}
                  \item Lorem ipsum dolor sit amet, consectetur adipiscing elit
                  \item Nam cursus est eget velit posuere pellentesque
              \end{itemize}
        \item Nulla commodo, erat quis gravida posuere, elit lacus lobortis est, quis porttitor odio mauris at libero
    \end{itemize}

    \bigskip % Vertical whitespace

    \begin{enumerate}
        \item Nam cursus est eget velit posuere pellentesque
        \item Vestibulum faucibus velit a augue condimentum quis convallis nulla gravida
    \end{enumerate}
\end{frame}

%------------------------------------------------

\section{Table and Figure Examples}


%------------------------------------------------

\subsection{Figure}

\begin{frame}
    \frametitle{Figure}

    \begin{figure}
        \includegraphics[width=0.8\linewidth]{creodocs_logo.pdf}
        \caption{Creodocs logo.}
    \end{figure}
\end{frame}
\subsection{Figure}

%------------------------------------------------

\section{Mathematics}

\begin{frame}
    \frametitle{Definitions \& Examples}

    \begin{definition}
        A \alert{prime number} is a number that has exactly two divisors.
    \end{definition}

    \smallskip % Vertical whitespace

    \begin{example}
        \begin{itemize}
            \item 2 is prime (two divisors: 1 and 2).
            \item 3 is prime (two divisors: 1 and 3).
            \item 4 is not prime (\alert{three} divisors: 1, 2, and 4).
        \end{itemize}
    \end{example}

    \smallskip % Vertical whitespace

    You can also use the \texttt{theorem}, \texttt{lemma}, \texttt{proof} and \texttt{corollary} environments.
\end{frame}

%------------------------------------------------

\begin{frame}
    \frametitle{Theorem, Corollary \& Proof}

    \begin{theorem}[Mass--energy equivalence]
        $E = mc^2$
    \end{theorem}

    \begin{corollary}
        $x + y = y + x$
    \end{corollary}

    \begin{proof}
        $\omega + \phi = \epsilon$
    \end{proof}
\end{frame}

%------------------------------------------------

\begin{frame}
    \frametitle{Equation}

    \begin{equation}
        \cos^3 \theta =\frac{1}{4}\cos\theta+\frac{3}{4}\cos 3\theta
    \end{equation}
\end{frame}

%------------------------------------------------

\begin{frame}[fragile] % Need to use the fragile option when verbatim is used in the slide
    \frametitle{Verbatim}

    \begin{example}[Theorem Slide Code]
        \begin{verbatim}
			\begin{frame}
				\frametitle{Theorem}
				\begin{theorem}[Mass--energy equivalence]
					$E = mc^2$
				\end{theorem}
		\end{frame}\end{verbatim} % Must be on the same line
    \end{example}
\end{frame}

%------------------------------------------------

\begin{frame}
    Slide without title.
\end{frame}

%------------------------------------------------

\section{Referencing}

\begin{frame}
    \frametitle{Citing References}

    An example of the \texttt{\textbackslash cite} command to cite within the presentation:

    \bigskip % Vertical whitespace

    This statement requires citation \cite{p1,p2}.
\end{frame}

%------------------------------------------------

\begin{frame} % Use [allowframebreaks] to allow automatic splitting across slides if the content is too long
    \frametitle{References}

    \begin{thebibliography}{99} % Beamer does not support BibTeX so references must be inserted manually as below, you may need to use multiple columns and/or reduce the font size further if you have many references
        \footnotesize % Reduce the font size in the bibliography

        \bibitem[Smith, 2022]{p1}
        John Smith (2022)
        \newblock Publication title
        \newblock \emph{Journal Name} 12(3), 45 -- 678.

        \bibitem[Kennedy, 2023]{p2}
        Annabelle Kennedy (2023)
        \newblock Publication title
        \newblock \emph{Journal Name} 12(3), 45 -- 678.
    \end{thebibliography}
\end{frame}

%----------------------------------------------------------------------------------------
%	ACKNOWLEDGMENTS SLIDE
%----------------------------------------------------------------------------------------

\begin{frame}
    \frametitle{Acknowledgements}

    \begin{columns}[t] % The "c" option specifies centered vertical alignment while the "t" option is used for top vertical alignment
        \begin{column}{0.45\textwidth} % Left column width
            \textbf{Smith Lab}
            \begin{itemize}
                \item Alice Smith
                \item Devon Brown
            \end{itemize}
            \textbf{Cook Lab}
            \begin{itemize}
                \item Margaret
                \item Jennifer
                \item Yuan
            \end{itemize}
        \end{column}
        \begin{column}{0.5\textwidth} % Right column width
            \textbf{Funding}
            \begin{itemize}
                \item British Royal Navy
                \item Norwegian Government
            \end{itemize}
        \end{column}
    \end{columns}
\end{frame}

%----------------------------------------------------------------------------------------
%	CLOSING SLIDE
%----------------------------------------------------------------------------------------

\begin{frame}[plain] % The optional argument 'plain' hides the headline and footline
    \begin{center}
        {\Huge The End}

        \bigskip\bigskip % Vertical whitespace

        {\LARGE Questions? Comments?}
    \end{center}
\end{frame}

%----------------------------------------------------------------------------------------

\end{document}